\begin{trapbox}

	\begin{description}[style=nextline, leftmargin=2.8em, labelsep=0.5em, itemsep=0.35em]
		\item[\textbf{\textcolor{CTrap}{易错点一（忽略正整数）}}] 误以为 $n=0$ 时结论也成立，认为 $\binom{0}{0}=2^{-1}$。
		\item[\textbf{\textcolor{CTrap}{正确理解（正整数条件）：}}] 结论要求 $n\ge 1$。$n=0$ 时 $\binom{0}{0}=1$，而 $2^{n-1}=1/2$，结论不成立。
		\item[\textbf{\textcolor{CTrap}{错因分析（未考虑边界）：}}] 未注意 $n$ 为正整数，盲目推广到 $n=0$，忽略指数定义域。
	\end{description}

	\medskip
	\begin{description}[style=nextline, leftmargin=2.8em, labelsep=0.5em, itemsep=0.35em]
		\item[\textbf{\textcolor{CTrap}{易错点二（混淆指数）}}] 将结论错记为 $2^{\,n}$ 而非 $2^{\,n-1}$。
		\item[\textbf{\textcolor{CTrap}{正确理解（指数减一）：}}] 所有系数总和为 $2^{\,n}$，奇偶各占一半，故为 $2^{\,n-1}$。可由 $E=O$ 和 $E+O=2^{\,n}$ 推导。
		\item[\textbf{\textcolor{CTrap}{错因分析（记忆错误）：}}] 只记结论不记推导，忽略了“一半”的除法运算。
	\end{description}

	\medskip
	\begin{description}[style=nextline, leftmargin=2.8em, labelsep=0.5em, itemsep=0.35em]
		\item[\textbf{\textcolor{CTrap}{易错点三（奇偶不分）}}] 需用奇数项和时错用偶数项和的表达式，或认为两者不同。
		\item[\textbf{\textcolor{CTrap}{正确理解（奇偶相等）：}}] 偶数项和与奇数项和相等，均为 $2^{\,n-1}$，无需区分符号。
		\item[\textbf{\textcolor{CTrap}{错因分析（推导不明）：}}] 未理解 $E-O=0$ 的推导，误以为奇偶系数和有差异。
	\end{description}

\end{trapbox}
